% §III --- System Model. Status: DRAFT.
\section{BUBBLES-Conformant System Model}\label{sec:system}

We consider $N$ small UAVs operating in a shared low-altitude volume and
answering a stream of VQA queries through an edge server over an A2G link.
Rather than invent the airspace, mobility, and conflict models, we anchor
every parameter to the SESAR JU BUBBLES D2.1
deliverable~\cite{bubbles2022d21}, double-citing peer-reviewed or regulatory
sources for load-bearing claims~\cite{weinert2022well,eu2021664}.

\subsection{Airspace, envelope, and demand}
The operating volume is a $5\times5$\,km area with three flight layers at
$20$--$120$\,m altitude (BUBBLES Appendix~G). The traffic
\emph{performance envelope} (Appendix~B) fixes the UAV kinematics: cruise
speed $5$--$25$\,m/s, climb rate $3$--$5$\,m/s, and airframe size
$0.5$--$5$\,m. The Appendix-G Table~G-13 operational demand---daily-mean
$8.17$ concurrent aircraft, peak $20$, $784$ operations/day, mean operation
duration $15$\,min---defines two operating conditions we evaluate under:

\begin{itemize}
\item \textbf{peak (\texttt{utm\_conflict})}: $20$ concurrent aircraft,
dense arrivals; the conflict constraint is genuinely activated (a stress
test of safety behaviour);
\item \textbf{nominal}: $8.17$ concurrent aircraft; the default evaluation
condition, where the airspace is uncongested.
\end{itemize}

Each aircraft is placed on an \emph{Area4D} operational
volume~\cite{bubbles2022d21}---a disk of radius $130$\,m matching the
nominal sensor footprint---on a jittered grid, and its background
operational intents are generated from that aircraft's own task stream
through a density gate (activation probabilities $0.30$ nominal /
$0.12$ peak). \emph{Honesty note}: we therefore do \emph{not} inject an
independent environmental traffic population; the demand is a
\emph{fleet-internal} traffic model in which every intent traces to a
served UAV, and a separate exogenous-traffic layer is left to future work.
A consequence of the Area4D packing is a scenario-intrinsic baseline
geometric-conflict rate---it is a property of the demand geometry, not an
artifact of the scheduling policy.

VQA tasks arrive as a VisDrone-derived query stream~\cite{zhu2021visdrone},
modelled as an operations-to-queries process in which each active operation
emits a question stream over its dwell. Each episode serves $20$ queries in
one-second decision slots (one service decision per slot); under the peak
preset the queries arrive as a front-loaded burst with a $50\%$
critical/high risk mix, while nominal arrivals are staggered.

\subsection{A2G link, bandwidth, and interference}\label{subsec:a2g}
The air-to-ground uplink uses log-distance path loss with an
elevation-dependent line-of-sight (LoS) term: we adopt the Al-Hourani LoS
sigmoid~\cite{alhourani2014optimal} with its original urban parameters
($a{=}9.61$, $b{=}0.16$) and a path-loss exponent of $2.2$--$2.25$,
cross-checked against 3GPP TR~36.777 aerial
channel models~\cite{3gpp2020tr36777}. Shadow fading uses
$\sigma_{\mathrm{SF}}=2$\,dB (36.777 UMa-AV LoS gives
$4.64\,e^{-0.0066h}\!\approx\!2.6$\,dB at $90$\,m), and small-scale fading is
Rician with $K{=}6$\,dB, within the $K\!\approx\!2$--$15$\,dB range reported
for measured A2G links~\cite{khawaja2019survey}. The received SNR is
quantized to $\{-5,\dots,20\}$\,dB bins.

\emph{Bandwidth and multiple access.} The A2G bandwidth is not a private
per-aircraft channel: it is a \emph{share of a common spectrum pool}, and
the continuous resource action allocates a fraction of that pool
($0.65$--$1.0$\,MHz per served task in our scenario tiers). The main
training and evaluation loop is \emph{slotted-serial on the uplink}: each
decision slot serves a single task upload (an implementation fact), so
concurrent co-channel interference does not accumulate on the main path and
the noise floor is the aggregate out-of-system interference power constant
$I_0$. We sweep $I_0$ from $-120$ to $-116$\,dBm as an
\emph{interference-floor axis} for zero-shot robustness; this is a raised
noise floor, not a claim of resolved multi-user contention. A partially
overlapping shared-pool leakage mechanism (concurrent uploads contributing
weighted co-channel interference) is exercised only in the multi-UAV
concurrency benchmark of Section~\ref{sec:experiments} and is left as an
extension of the main serial model. Transmission delay uses a
capacity-abstraction rate (Shannon spectral efficiency at the operating
SNR, with an implementation efficiency gap of $\eta\!\approx\!0.75$
understood but not separately modelled), so quoted uplink times are
lower bounds on a fielded modem.

\subsection{Tactical-contact separation chain}
Following BUBBLES §3.3.4 (Block~4), the tactical-contact separation minimum
is built up from the physical near-mid-air-collision distance by adding
error and time-window terms:
\begin{equation}\label{eq:chain}
\begin{split}
\dtc = {}&\dnmac + \mathrm{TSE} + \mathrm{SSE} \\
       &+ V_c\big[(T_1{+}T_2) + T_{\mathrm{res}} + (T_3{+}T_4)\big],
\end{split}
\end{equation}
where $\dnmac$ is the downstream near-collision barrier, TSE/SSE are the
total/surveillance system errors, $V_c$ is the closing speed, $T_1,T_2,T_3$
are detection/decision/manoeuvre time windows, and $T_{\mathrm{res}}$ is
the separation-recovery time. The final term $T_4$ is the sum of three
normally distributed components---tracking delay, \emph{separation
communication delay}, and pilot-execution delay (Table~G-2 example:
communication mean $1.8$\,s / $\sigma\,1.0$, $T_4$ total mean $7.82$\,s /
$\sigma\,4.24$). Equation~\eqref{eq:chain} is the mechanistic core of our
safety bridge (Section~\ref{sec:experiments}): communication
delay~$\uparrow$ $\Rightarrow$ $T_4\uparrow$ $\Rightarrow$
$V_cT_4\uparrow$ $\Rightarrow$ $\dtc\uparrow$ $\Rightarrow$ airspace
capacity~$\downarrow$. Communication delay is not exogenous noise; it is an
additive Gaussian component inside the separation minimum.

A pair of UAVs is in \emph{tactical conflict} only when both a spatial and
a temporal criterion hold: the closest-point-of-approach (CPA) distance
falls below $\dtc$ \emph{and} time-to-CPA is below $60$\,s. Critically, CPA
is extrapolated from the kinematics of the \emph{assigned} UAV (the one
tasked to serve), not the nearest arbitrary UAV; using the nearest UAV
would make conflicts almost exogenous to the policy, leaving ``select
cache'' as the only avoidance action (see Section~\ref{sec:algorithm}).

\subsection{How payload enters $T_4$: an M/G/1 queueing coupling}
\label{subsec:mg1}
The separation-communication component of $T_4$ (Table~G-2 mean $1.8$\,s)
is the tactical loop's own exchange---formulate the resolution, deliver the
instruction. Our VQA payload does not \emph{replace} that exchange; it
\emph{contends} with it on the shared RAN. We therefore treat the $1.8$\,s
as the queue-empty tactical baseline and add the contention as the mean
waiting time $W$ of a non-preemptive priority M/G/1
queue~\cite{kleinrock1975queueing}: the separation-relevant (C2-class)
message is the high-priority class, payload uploads are the low-priority
class whose service time $S$ is the airtime of the evidence level the
scheduler chose, and per-class utilisations follow the G-13 demand
($\rho_{\mathrm{C2}}/\rho_{\mathrm{payload}}=0.10/0.35$ nominal,
$0.25/0.55$ peak). Then
\begin{equation}\label{eq:mg1}
T_4^{\mathrm{comm}} = 1.8\,\mathrm{s} + W,\qquad
W = W_{\mathrm{M/G/1}}\!\big(S(\text{evidence}),\,\rho\big),
\end{equation}
so a heavy upload lengthens the \emph{residual service} a tactical message
must wait behind, and the evidence choice reaches $\dtc$ through the queue
rather than by substitution. Because whether C2 shares the payload link is
contested (Section~\ref{subsec:qos}), we evaluate both a \emph{shared}
mode (C2 in the pool as the top class) and a \emph{dedicated} mode (C2 on
its own band; only same-class payload contention remains).
Table~\ref{tab:wsens} shows $W$ at the $-5$\,dB worst case: a token upload
($14$\,ms airtime) adds $3$--$8$\,ms of wait---the tactical loop does not
notice---while an image upload ($3.72$\,s airtime) adds $0.78$--$1.98$\,s,
a $44$--$110\%$ inflation of the separation-communication mean.

\begin{table}[t]
\centering
\caption{M/G/1 wait $W$ (s) added to the $1.8$\,s separation-comm baseline
at $-5$\,dB Rayleigh, by load and C2 spectrum mode.}
\label{tab:wsens}
\footnotesize
\begin{tabular}{@{}lcccc@{}}
\toprule
 & \multicolumn{2}{c}{shared C2} & \multicolumn{2}{c}{dedicated C2} \\
\cmidrule(lr){2-3}\cmidrule(lr){4-5}
Evidence & nominal & peak & nominal & peak \\
\midrule
$s_1$ token ($S{=}0.014$\,s) & $0.004$ & $0.008$ & $0.003$ & $0.005$ \\
$s_2$ image ($S{=}3.72$\,s)  & $0.930$ & $1.984$ & $0.783$ & $1.435$ \\
\bottomrule
\end{tabular}
\end{table}

\subsection{Deadline QoS and the C2 priority class}\label{subsec:qos}
Each query carries a delay budget $\tau_k$ interpreted as the
\emph{communication-and-decision window}: the time from question arrival to a
returned answer usable by the separation service. Flight-repositioning time
is a \emph{tasking-layer} quantity and is deliberately \emph{not} charged to
$\tau_k$. BUBBLES gives no hard communication-latency numbers---a
fit-for-purpose principle citing ED-261 GEN-SUR (D2.1 p.\,62)---so rather
than assert a deadline we anchor it to the separator-communication delay of
the tactical chain: Table~G-2 models that delay as $\mathcal{N}(1.8,
\sigma\,1.0)$\,s, and we set the critical and normal budgets at its
$+2\sigma$/$+3\sigma$ confidence points, $\tau_{\mathrm{critical}}=3.8$\,s
and $\tau_{\mathrm{normal}}=4.8$\,s. Treating the separation-relevant
semantic evidence stream as the highest-priority class within the
task-data-link QoS is consistent with the C2 key-performance-indicator
targets of 3GPP TS~22.125~\cite{3gpp2022ts22125} and with EUROCAE ED-282 as
supporting guidance~\cite{eurocae2020ed282}. We stress the scope red line:
authenticated command-and-control / control-and-non-payload-communication
(CNPC) runs on protected aeronautical spectrum (ICAO WRC-12
$5030$--$5091$\,MHz) and is \emph{outside} this paper; a structural
precedent for isolating such a flow as a separate QoS flow on the same RAN
is 3GPP TS~22.125 / TS~23.256~\cite{3gpp2022ts22125,3gpp2022ts23256}.

\subsection{Spec-attainability certificate and the escalation budget}
\label{subsec:cert}
Peak overload creates queries that \emph{no} action can serve within spec:
at a weak link even the best transmission service cannot clear the quality
threshold inside the deadline. Charging such physically unserveable
queries to the policy as quality violations would poison both the learning
signal and the safety accounting, so we make unserveability an explicit,
certified state.

\begin{definition}[Spec-attainability certificate]\label{def:cert}
At slot $t$, query $k$ with remaining window $\tilde\tau_k(t)$ and quality
threshold $\varepsilon_k$ is \emph{spec-attainable},
$\mathrm{A}_k(t)=1$, iff some transmission service $s\in\{s_1,s_2\}$
satisfies both
$\mathrm{LCB}\big(q_s(\gamma_k)\big)\ge\varepsilon_k$ and
$D^{\mathrm{comm}}_s(\gamma_k)\le\tilde\tau_k(t)$,
where $q_s$ is the calibrated quality of service $s$ at link state
$\gamma_k$, LCB its lower confidence bound, and $D^{\mathrm{comm}}_s$ the
communication-window delay of Section~\ref{subsec:qos}.
\end{definition}

The certificate is an \emph{environment} property (calibrated quality
tables plus link physics, independent of the learner), is exposed to the
policy as a state feature, and is recomputed \emph{every} slot from the
remaining window---a query whose window has been eaten by deferrals turns
unattainable mid-episode. It drives a certified \emph{escalation} path for
the safety-critical class: when a critical/high query ends the episode
rejected or expired, the certificate decides the verdict. If
$\mathrm{A}_k=1$ the policy dropped a serveable safety query and is charged
the full quality violation; if $\mathrm{A}_k=0$ the query is
\emph{escalated}---handed to a higher-tier service (relay tasking, human
review) outside this scheduler---carrying no quality cost but consuming a
budgeted escalation event. Escalation is priced by a dedicated dual channel
with budget $\delta_{\mathrm{esc}}$ (Section~\ref{sec:problem}).

The budget is \emph{derived, not tuned}: rolling the certificate over the
scenario arrival mix measures the structurally unserveable fraction of the
critical load ($0.105$ under peak, $0.000$ under nominal), and
$\delta_{\mathrm{esc}}$ adds a $0.05$ headroom:
$\delta_{\mathrm{esc}}^{\mathrm{peak}}=0.155$,
$\delta_{\mathrm{esc}}^{\mathrm{nominal}}=0.05$. Queries that are not
escalated form the \emph{admitted set}; all rate constraints in
Section~\ref{sec:problem} are stated on it. This is the standard
demanded-versus-achievable split of admission control: the scheduler is
accountable for what it admits, and the certificate makes the
un-admittable load visible, budgeted, and auditable rather than silently
absorbed into violation statistics.

\subsection{TLS anchoring and the conflict budget}
The target level of safety is a total $\mathrm{TLS}=10^{-6}$ fatal
accidents per flight hour (FAT/FH), of which the mid-air-collision share is
$\mathrm{TLS}_{\mathrm{MAC}}=2.5\times10^{-7}$ FAT/FH; collision risk scales
as $N(N{-}1)$ under the gas-kinetic model, and a severity ladder
MAC$\leftarrow$NMAC$\leftarrow$IC$\leftarrow$SL$\leftarrow$TC places one
mitigation barrier at each level. We use this ladder to \emph{derive} a
per-slot conflict budget rather than assert one: a tactical-contact
violation sits at the top of the ladder, protected by several downstream
barriers, so it may occur far more frequently than a MAC. Normalizing that
admissible frequency by the peak concurrency and the slot window yields a
per-slot tactical-conflict budget of $\approx\!0.08$, which we use as the
conflict constraint limit in Section~\ref{sec:problem}.
\emph{Honesty note}: BUBBLES gives no hard communication
reliability/latency numbers (a fit-for-purpose principle), so the specific
normalization constant behind $0.08$ is our engineering design choice
within the barrier-ladder$+$demand framework, declared as such rather than
read directly from the deliverable.

\subsection{Separation kernel and validation gate}
The chain~\eqref{eq:chain}, the TLS anchoring, and the CPA dual criterion
are implemented in a standalone separation kernel gated by a reproduction
test against the official numerics: it recomputes BUBBLES Tables~G-4/G-5
cell-by-cell to within $0.18\%$ (e.g., a total contact distance of
$370.53\to370.49$\,m for a SAIL~I--II scenario), and a smoke scenario
injecting an approaching pair lifts the measured conflict rate from $0$ to
$0.5$, confirming the dual criterion fires; the full suite passes $14/14$
unit tests. The kernel is a fixed ``consumer'' that receives the measured
link delay from the semantic layer and returns separation distances and
conflict flags.

\subsection{Semantic evidence levels}
The transmitter chooses, per query, one of three evidence levels (a fourth,
region-of-interest crop, is retained as an optional action but excluded
from the main comparison); the receiver is a frozen commercial
vision--language model (Qwen2-VL-2B)~\cite{wang2024qwen2vl}. This taxonomy
and its per-question complementarity are established in the companion
paper~\cite{zhang2026askb4tx}; we summarize the transmission cost, since it
drives both the resource and safety constraints.
\begin{itemize}
\item $s_0$ \textbf{cache}: reuse a fresh cached answer; $0$\,B, no
transmission, and---in the BUBBLES setting---\emph{zero conflict
exposure}, since no flight/communication action is initiated. A cache entry
ages fresh$\to$stale$\to$expired and cannot be reused past expiry; its
effective accuracy is $A_{s_0}=A_{\mathrm{LUT}}(0.85+0.15\,p_{\mathrm{hit}})$.
\item $s_1$ \textbf{token}: light detector tokens (class/box/confidence),
$\sim\!0.7$--$1.2$\,KB, framed under LDPC rate-$1/2$ with
BPSK~\cite{gallager1962ldpc}, costing $N_{s_1}=16\,B_{\mathrm{tok}}$ complex
channel uses ($\sim\!11$k--$19$k), decoded by a symbolic reader or the VLM.
\item $s_2$ \textbf{image}: full JPEG at rate-adaptive quality, $\sim\!110$\,KB,
costing $N_{s_2}=8B_{\mathrm{JPEG}}/\mathrm{SE}(\bar\gamma)$ uses; at $-5$\,dB
this reaches $3.72$M uses ($\approx\!265\times$ the token cost) and a
$3.72$\,s uplink, answered by the VLM.
\end{itemize}
The two-order-of-magnitude gap between $s_1$ and $s_2$ under weak links is
what makes evidence selection a scheduling lever with both a
resource-budget and a flight-safety effect.
